% carleson_edgeband.tex -- round 482.
% PRIME.RDAGGER.CARLESON_EDGEBAND.01
% NO RH CLAIM. NO anti-RH CLAIM.
\documentclass[11pt]{article}

\usepackage[a4paper,margin=2.1cm]{geometry}
\usepackage{amsmath,amssymb,amsthm}
\usepackage{booktabs,array}
\usepackage{microtype}
\usepackage{xcolor}
\usepackage[colorlinks=true,linkcolor=blue!50!black,
            urlcolor=blue!50!black]{hyperref}

\newtheorem{lemma}{Lemma}
\newtheorem{observation}{Observation}
\theoremstyle{remark}
\newtheorem{remark}{Remark}
\newcommand{\code}[1]{\nolinkurl{#1}}
\newcommand{\qdag}{q^\dagger}
\newcommand{\scr}{\mathrm{SCR}}
\newcommand{\arch}{\mathrm{ARCH}}
\setlength{\emergencystretch}{2em}

\title{\bfseries Carleson anatomy at the spectral edge\\[2mm]
\large PRIME.RDAGGER.CARLESON\_EDGEBAND.01 --- round 482}
\author{Stefan Hamann}
\date{August 31, 2026}

\begin{document}
\maketitle

\begin{abstract}\noindent
The r460 spectral measure is measured on every dyadic scale
$2^{-1},\ldots,2^{-20}$ for the interval-certified selected windows
$k=5,\ldots,16$.  A finite uniform envelope exists:
\[
 F_k(t)\le0.666528\,t^{1.25}\qquad(0<t\le\tfrac12,\quad5\le k\le16),
\]
where the constant is taken over every spectral jump, not just the
dyadic samples.  This is not a cofinal theorem.  More decisively, even
the separately optimized Carleson envelope of each exact finite spectrum
gives a sufficient-rule value between $1.055000$ and $2.700019$:
it closes none of the twelve windows whose direct $\qdag_k$ is already
below one.  Christoffel and local-Plancherel majorants lose the measured
edge alignment; the only scramble-sensitive large-sieve formulation is
the target projector inequality itself, with no classical transfer
through the data-dependent orthogonal-polynomial basis.  The verdict is
\[
\boxed{\begin{gathered}
\text{\code{STUCK(SR_FAILS_FINITE_CENSUS;}}\\
\text{\code{RESTRICTED_LARGE_SIEVE_IS_TARGET;}}\\
\text{\code{STRICT_F_SCRAMBLE_PREMISE_FALSE)}}.
\end{gathered}}
\]
No RH claim and no anti-RH claim.
\end{abstract}

\section{Object and the sufficient rule}

For the production selected window
\[
 a_k=2^k,\quad r_k=\lfloor\sqrt{k}\rfloor,\quad
 m_k=k2^{r_k}-1,\quad \Delta_k=2^{-r_k}\log2,
\]
let $C_k$ be the r460 negative-sampling contraction and
$u_k=b_k/\sqrt{B_k}$ the budget-normalized border vector in the
$\mu_k$-orthonormal polynomial basis.  If
$C_ke_{k,j}=\lambda_{k,j}e_{k,j}$, put
\[
 d_{k,j}=1-\lambda_{k,j},\qquad
 w_{k,j}=|\langle u_k,e_{k,j}\rangle|^2,\qquad
 F_k(t)=\sum_{0<d_{k,j}\le t}w_{k,j}.
\]
Thus $\qdag_k=\sum_jw_{k,j}/d_{k,j}
=\int(1-\lambda)^{-1}\,d\sigma_k(\lambda)$.
All MAIN eigenvalues in this census lie below one.

\begin{lemma}[Carleson sufficient rule]\label{lem:sr}
If $F_k(t)\le Ct^{1+\varepsilon}$ for $0<t\le\eta$, then
\[
 \qdag_k\le
 \frac{\sigma_k([0,1-\eta])}{\eta}
 +C(1+\varepsilon^{-1})\eta^\varepsilon.
 \tag{SR}
\]
In particular, a strict upper bound below one implies $\qdag_k<1$.
\end{lemma}

\begin{proof}
With $x=1-\lambda$ and $F(t)=\sigma(\{0<x\le t\})$,
Stieltjes integration by parts gives
\[
 \int_{(0,\eta]}\frac{dF(x)}x
 =\frac{F(\eta)}{\eta}+\int_0^\eta\frac{F(x)}{x^2}\,dx
 \le C\eta^\varepsilon+\frac C\varepsilon\eta^\varepsilon.
\]
On $x\ge\eta$, $x^{-1}\le\eta^{-1}$.  At an atom on the splitting
boundary the displayed closed intervals may double-count; the
inequality remains valid.  The computation below minimizes between
adjacent atoms, so it does not rely on equality at a boundary.
\end{proof}

\section{Dyadic anatomy}

Tables~\ref{tab:f1} and~\ref{tab:f2} give the full requested curve.
Each entry is $F_k(2^{-j})$.  Zeros are actual empty edge bands, not
rounded positive values.

\begin{table}[ht]
\centering\scriptsize
\caption{Dyadic edge mass, scales $j=1,\ldots,10$.}
\label{tab:f1}
\begin{tabular}{@{}r*{10}{r}@{}}
\toprule
$k\backslash j$&1&2&3&4&5&6&7&8&9&10\\
\midrule
5&7.95e-3&7.95e-3&0&0&0&0&0&0&0&0\\
6&1.05e-2&1.75e-6&1.75e-6&1.75e-6&0&0&0&0&0&0\\
7&2.12e-3&1.18e-3&5.55e-4&5.55e-4&5.55e-4&5.55e-4&0&0&0&0\\
8&8.46e-2&1.82e-3&1.56e-3&2.76e-4&2.76e-4&2.76e-4&2.76e-4&0&0&0\\
9&9.03e-2&5.38e-2&1.45e-3&6.94e-5&6.94e-5&3.23e-5&1.22e-5&1.22e-5&1.22e-5&1.22e-5\\
10&7.78e-2&9.45e-3&7.88e-4&3.72e-4&8.35e-5&4.90e-5&4.90e-5&4.90e-5&2.74e-6&2.74e-6\\
11&1.66e-1&1.79e-3&1.67e-3&5.22e-5&3.81e-5&3.81e-5&9.85e-6&9.85e-6&9.85e-6&2.44e-6\\
12&1.55e-1&4.65e-2&1.45e-3&1.22e-3&3.43e-5&2.51e-5&1.19e-5&9.16e-6&9.16e-6&9.16e-6\\
13&1.01e-1&4.69e-2&8.15e-4&4.75e-4&1.07e-4&1.68e-5&1.31e-5&5.14e-6&5.14e-6&3.26e-6\\
14&1.20e-1&8.05e-2&5.65e-4&3.87e-4&3.50e-4&1.08e-5&1.08e-5&4.25e-6&3.28e-6&3.28e-6\\
15&1.03e-1&6.48e-2&1.84e-2&2.71e-4&1.67e-4&1.16e-4&1.16e-5&1.14e-5&3.43e-6&1.95e-6\\
16&1.36e-1&1.11e-2&6.41e-4&4.00e-4&6.14e-5&2.22e-5&8.37e-6&4.39e-6&7.03e-7&4.82e-7\\
\bottomrule
\end{tabular}
\end{table}

\begin{table}[ht]
\centering\scriptsize
\caption{Dyadic edge mass, scales $j=11,\ldots,20$.}
\label{tab:f2}
\begin{tabular}{@{}r*{10}{r}@{}}
\toprule
$k\backslash j$&11&12&13&14&15&16&17&18&19&20\\
\midrule
5&0&0&0&0&0&0&0&0&0&0\\
6&0&0&0&0&0&0&0&0&0&0\\
7&0&0&0&0&0&0&0&0&0&0\\
8&0&0&0&0&0&0&0&0&0&0\\
9&0&0&0&0&0&0&0&0&0&0\\
10&2.74e-6&0&0&0&0&0&0&0&0&0\\
11&2.44e-6&2.44e-6&0&0&0&0&0&0&0&0\\
12&6.73e-7&6.73e-7&6.73e-7&0&0&0&0&0&0&0\\
13&3.26e-6&4.73e-7&4.73e-7&4.73e-7&0&0&0&0&0&0\\
14&1.29e-6&1.29e-6&2.45e-7&2.45e-7&2.45e-7&0&0&0&0&0\\
15&1.95e-6&6.91e-7&6.91e-7&1.21e-7&1.21e-7&1.21e-7&0&0&0&0\\
16&4.07e-7&3.23e-7&5.18e-8&4.96e-8&4.96e-8&7.27e-9&7.27e-9&0&0&0\\
\bottomrule
\end{tabular}
\end{table}

\subsection{Fits, AIC, and drops}

The regression uses only positive dyadic values:
$\log F=\log C+(1+\varepsilon)\log t$.  ``drop'' is the disclosed
number of zero dyads.  $\Delta\mathrm{AIC}$ compares the free exponent
(two parameters) to exponent one (one parameter); negative favours the
free exponent.  This regression is descriptive.  The last column is
instead the rigorous envelope constant for the frozen spectrum,
\[
 C_k^{(.25)}=\max_{\substack{d_{k,j}\le1/2\\m\le j}}
 \frac{\sum_{i\le m}w_{k,i}}{d_{k,m}^{1.25}},
\]
taken at every exact jump after ordering the positive gaps.

\begin{table}[ht]
\centering\small
\caption{Power fits and exact finite-spectrum envelopes.}
\label{tab:fits}
\begin{tabular}{@{}rrrrrrr@{}}
\toprule
$k$&$n$&drop&$C_{\rm reg}$&$\varepsilon_{\rm reg}$&
$\Delta$AIC&$C_k^{(.25)}$\\
\midrule
5&2&18&---&---&---&0.071253\\
6&4&16&0.010469& 2.765281& -0.371&0.043504\\
7&6&14&0.001925&-0.630964&-10.397&0.123841\\
8&7&13&0.026820& 0.169318&  1.689&0.222095\\
9&10&10&0.048175&0.476566& -1.316&0.307378\\
10&11&9&0.035956&0.382989& -4.251&0.189380\\
11&12&8&0.012824&0.217173&  0.509&0.522558\\
12&13&7&0.052105&0.417614& -4.372&0.523342\\
13&14&6&0.027714&0.299983& -2.768&0.327817\\
14&15&5&0.027108&0.287465& -2.409&0.666528\\
15&16&4&0.051598&0.329506& -6.315&0.620111\\
16&17&3&0.018161&0.360118&-10.640&0.327005\\
\bottomrule
\end{tabular}
\end{table}

Pooling all $127$ positive observations gives
$C_{\rm reg}=0.0190020$, $\varepsilon_{\rm reg}=0.274663$,
AIC $=137.991$ and $\Delta\mathrm{AIC}=-23.696$.  This pooled slope
does not remove the cross-window instability: the individual exponent
excess ranges from $-0.631$ to $+2.765$.  On the exact finite spectra,
not merely the dyadic grid, the uniform pair
\[
 (C,\varepsilon,\eta)=(0.666527842,\;0.25,\;0.5)
\]
works for all $k=5,\ldots,16$.  Finiteness makes existence of some
envelope unsurprising; this row says nothing about infinitely many $k$.

\section{The SR balance}

For each $k$ and each $\varepsilon\in(0,4]$ the code orders all exact
spectral gaps.  Between two consecutive gaps the edge constant is fixed,
so $A/\eta+B\eta^\varepsilon$ is minimized analytically at its stationary
point or an interval endpoint.  A deterministic grid locates the
$\varepsilon$ basin and a bounded scalar minimization refines it.

\begin{table}[ht]
\centering\small
\caption{Best separate finite-spectrum SR balance.  Margin is $1-$SR.}
\label{tab:sr}
\begin{tabular}{@{}rrrrrrrr@{}}
\toprule
$k$&$\eta$&$\varepsilon$&$C$&far&edge&SR&margin\\
\midrule
5 &1.000000&0.877030&0.214041&0.620952&0.458092&1.079044&-0.079044\\
6 &0.882140&2.473420&0.548160&0.490507&0.564493&1.055000&-0.055000\\
7 &0.617552&0.214229&0.106085&1.006840&0.542290&1.549130&-0.549130\\
8 &0.626667&0.311907&0.235612&0.735956&0.856583&1.592540&-0.592540\\
9 &0.999123&0.483865&0.425816&0.231699&1.305291&1.536990&-0.536990\\
10&0.677853&0.448722&0.218131&0.843784&0.591497&1.435280&-0.435280\\
11&0.286163&0.224233&0.098080&1.597035&0.404484&2.001520&-1.001520\\
12&0.375508&0.410194&0.356707&0.970016&0.820565&1.790581&-0.790581\\
13&0.789145&0.361612&0.381615&0.384508&1.319005&1.703513&-0.703513\\
14&0.154276&0.103930&0.029049&2.445934&0.254085&2.700019&-1.700019\\
15&0.512828&0.410749&0.838622&0.373879&2.189326&2.563205&-1.563205\\
16&0.972894&0.441995&0.416970&0.158621&1.343929&1.502550&-0.502550\\
\bottomrule
\end{tabular}
\end{table}

Every direct value is live,
$\qdag_k\in[0.852160,0.945017]$, but every optimized SR majorant is
above one.  Thus the loss is in replacing a many-step spectral
distribution by one power envelope, not in the underlying finite
positivity.  The r460 two-band calculation and the present calculation
share the same failure: a coarse far-field denominator and a separate
edge majorant discard the detailed placement of the atoms.

\section{World controls: strict edge versus dangerous edge}

The reviewer premise needs one correction.  The r467 factor
$28.5$--$5.3\cdot10^5$ concerned the normalized overlap with the
dangerous top eigenvector.  It was not a measurement of
$F^\scr_k(t)=\sigma^\scr_k([1-t,1))$.  SCRAMBLE and ARCH generally
have eigenvalues at or above one, which strict $F$ excludes.  Define
the control diagnostic
\[
 D_k(t)=\sigma_k([1-t,\infty)).
\]
It agrees with $F_k(t)$ in the MAIN world but also records a
supercritical instability.  At $t=1/2$, the normalized fractions
$D/\|u\|^2$ are as follows.

\begin{table}[ht]
\centering\small
\caption{Strict masses and normalized dangerous masses at $t=1/2$.}
\label{tab:world}
\begin{tabular}{@{}rrrrrrr@{}}
\toprule
$k$&$F^{\rm M}$&$F^\scr$&$F^\arch$&
$D^{\rm M}/\gamma_{\rm M}$&$D^\scr/\gamma_\scr$&
$D^\arch/\gamma_\arch$\\
\midrule
5 &7.95e-3&3.12e0&0&0.0095&0.2985&0.9993\\
6 &1.05e-2&0&0&0.0133&0.2297&1.0000\\
7 &2.12e-3&0&0&0.0034&0.0503&1.0000\\
8 &8.46e-2&5.74e3&8.06e-3&0.1551&0.4992&1.0000\\
9 &9.03e-2&2.52e4&0&0.1436&0.8342&1.0000\\
10&7.78e-2&1.25e0&2.32e-5&0.1181&0.9579&1.0000\\
11&1.66e-1&2.67e-1&0&0.3477&0.9957&1.0000\\
12&1.55e-1&4.18e0&0&0.3558&0.9766&1.0000\\
13&1.01e-1&1.61e-1&0&0.1925&0.9978&1.0000\\
14&1.20e-1&1.16e0&0&0.3182&0.9999&1.0000\\
15&1.03e-1&3.05e0&9.19e-5&0.2953&0.9839&1.0000\\
16&1.36e-1&7.85e-1&7.85e-6&0.2444&1.0000&1.0000\\
\bottomrule
\end{tabular}
\end{table}

The dangerous fraction is larger under SCRAMBLE for every $k$, and
ARCH is at least $0.9993$.  The world test therefore confirms the r467
joint-alignment diagnosis.  But strict $F^\scr(1/2)<F^{\rm M}(1/2)$
at $k=6,7$: its mass moved through $\lambda=1$.  A proof mechanism
must be sensitive to this distinction; treating zero strict mass as a
good scrambled control is a category error.

\section{The three proposed mechanisms}

\subsection{Christoffel functions}

Let $\mathcal P_{<N_k}$ carry the $L^2(\mu_k)$ norm and let
$\ell_k$ be the normalized border functional represented by $u_k$.
The Christoffel/RKHS inequality is
\[
 |\ell_k(p)|^2\le\|\ell_k\|_{\mathcal P_{<N_k}^*}^2
 \|p\|_{\mu_k}^2=\|u_k\|^2\|p\|_{\mu_k}^2.
\tag{Chr}
\]
Pointwise Christoffel bounds followed by a triangle inequality over
the border atoms are weaker.  Applied to the top unit eigenvector,
(Chr) exceeds the measured squared overlap by factors
$105$ through $7.64\cdot10^7$.  More importantly it contains no $t$.
Classical endpoint asymptotics for regular fixed measures do not apply
uniformly to this moving, discrete, source-dependent $\mu_k$; proving
the needed regularity is the retained orthogonal-polynomial/RH problem.

\paragraph{Scramble and alias.}
The numerical Christoffel kernel changes under scramble only because
the entire measure changes; the abstract bound does not retain the
arithmetic node placement and gives no favourable direction.  This is
r469 candidate (g), the Loewner/RKHS majorant, and it again separates
the extremal data from its edge mode.  Verdict: \textbf{KILL}.

\subsection{Discrete large sieve}

Put $E_k(t)=\mathbf1_{[1-t,1)}(C_k)\mathbb R^{N_k}$.  The exact
restricted inequality required is
\[
 \sum_{e_{k,j}\in E_k(t)}
 |\ell_k(e_{k,j})|^2
 =\|E_k(t)u_k\|^2
 \le C t^{1+\varepsilon}.
\tag{RLS}
\]
The literal arithmetic positions enter twice:
\[
 \{\log n\}\longrightarrow c_k^P\longrightarrow(\mu_k,\nu_k)
 \longrightarrow\{p_{k,i}\}\longrightarrow(C_k,u_k).
\]
This joint dependence is exactly why (RLS) passes the methodological
scramble filter.  Numerically, at $\varepsilon=0.25$, the sharp finite
constants on $t\le1/2$ are $0.0435$--$0.6665$ (apart from $k=5$,
$0.0713$); the factor to the binding measurement is one by construction.
That is a calibration, not a proof.

A classical discrete large sieve controls an exponential sum in terms
of node separation and coefficient energy \emph{before} this nonlinear
Stieltjes/orthogonalization map.  No known theorem transfers such a
bound through the MAIN-dependent basis while preserving a positive
power of the singular-value defect $t$.  Taking the supremum over unit
$p\in E_k(t)$ shows that (RLS) is exactly the Carleson target, not a
strict reduction of it.  Replacing it by a raw-node large sieve repeats
the r467 source-to-border-basis gap.  Verdict: \textbf{BEST TYPING,
BUT TARGET-EQUIVALENT; OPEN}.

\subsection{Local Plancherel}

Let $Q_k$ project onto the lowest quarter of the orthonormal DCT modes
in coefficient degree and let $P_k(t)$ be the edge projector.  The
exact elementary split
\[
 \|P_k(t)u_k\|
 \le\|Q_ku_k\|+\|P_k(t)(I-Q_k)\|\,\|u_k\|
\tag{Pl}
\]
is a local-Plancherel budget.  At $t=2^{-8}$ and $k=9,\ldots,16$ its
square exceeds the measured $F_k(t)$ by factors
$5.52\cdot10^3$--$3.59\cdot10^4$.  The low-frequency character of
individual top modes does not imply a small operator leakage for the
whole edge subspace, and the triangle split loses their phase relation
with $u_k$.

\paragraph{Scramble and alias.}
$P_k$ and $u_k$ both change under scramble, but inequality (Pl) treats
their two contributions separately.  It is the finite-coordinate form
of r469 candidate (c), whose missing edge symbol/Riemann--Hilbert
asymptotics were already identified, and of the r467 Abel separation.
Verdict: \textbf{KILL AT THE $10^3$ GATE}.

\section{Mechanism ranking and closure status}

\begin{table}[ht]
\centering\small
\caption{Reviewer-candidate ranking.}
\begin{tabular}{@{}p{0.19\linewidth}p{0.18\linewidth}
                p{0.18\linewidth}p{0.34\linewidth}@{}}
\toprule
candidate&sharpness&scramble judgment&status / alias\\
\midrule
Christoffel&$105$--$7.64\cdot10^7$&abstract bound blind&
KILL; r469(g), separate-supremum majorant\\
restricted large sieve&factor $1$ tautologically&yes, literal nodes in
$C_k,u_k$&best typing; (RLS) equals target, raw-node version has r467 gap\\
local Plancherel&$5.52\cdot10^3$--$3.59\cdot10^4$&
objects change, estimate splits them&KILL; r469(c), r467 Abel class\\
\bottomrule
\end{tabular}
\end{table}

\noindent\fbox{\parbox{0.96\linewidth}{%
\textbf{Closing status, verbatim.}\quad
\code{STUCK(SR_FAILS_FINITE_CENSUS;}
\code{RESTRICTED_LARGE_SIEVE_IS_TARGET;}
\code{STRICT_F_SCRAMBLE_PREMISE_FALSE).}
The finite uniform Carleson envelope is real, but an independently
optimized SR envelope is $>1$ on every certified window.  The data do
not refute a uniform Carleson law on infinitely many selected $k$;
they do show that the measured constants and a single power envelope
do not recover even the known finite inequalities.  The only candidate
that retains the arithmetic/eigenbasis correlation is (RLS), which is
the desired projector estimate itself and has no applicable classical
literature theorem.  Hence neither \code{SR_PROVED} nor
\code{CARLESON_REDUCED} nor \code{CARLESON_REFUTED} is justified.}}

\section{Numerical sealing and kill table}

The probe reconstructs the r463-faithful active source $n<a_k$ and
reproduces the r459/r468 $\qdag$ pins.  The curve record has SHA-256
prefix \code{e45622d390abf932}.  For the frozen binary64 symmetric
matrix, if $E\widehat\Lambda E^T$ is the computed decomposition, Weyl's
theorem encloses every ordered eigenvalue within
$\|CE-E\widehat\Lambda\|_2\le2\cdot10^{-15}$.  The closest dyadic
threshold is more than $2.5\cdot10^8$ such radii away.  Independent
dps-70 diagonalization at $k=5,9,12,16$ changes the top eight
eigenvalues by at most $4.45\cdot10^{-16}$.

These are validated enclosures for the frozen matrix and certify every
dyadic membership decision.  They are not a fresh interval propagation
of all source logs through every eigenvector overlap; the rigorous
r465/r468 interval certificates remain the authority for $\qdag<1$.
Accordingly the $F$ ordinates are sealed high-precision census values,
not promoted interval theorems.

\begin{table}[ht]
\centering\small
\caption{Kill and honesty gates.}
\begin{tabular}{@{}lll@{}}
\toprule
gate&status&reason\\
\midrule
production object / $k=5..16$&PASS&r463 active source; q pins reproduced\\
dyadic curves / determinism&PASS&240 values; sealed curve hash\\
AIC / drops&PASS&zeros disclosed; free vs exponent-one comparison\\
finite uniform pair&PASS&exact jumps, not only dyads\\
SR $<1$&KILL&best $1.055000$; no finite closure\\
SCRAMBLE strict $F$&KILL as premise&two reversals from $\lambda\ge1$\\
SCRAMBLE dangerous mass&PASS&larger on all twelve worlds\\
Christoffel&KILL&no $t$; up to $7.64\cdot10^7$ loose\\
local Plancherel&KILL&minimum factor $5.52\cdot10^3$\\
large sieve&OPEN / TARGET&(RLS) is exactly the projector embedding\\
33-route alias check&PASS&31 dossier routes plus r467/r468 checked\\
edge dps / membership&PASS&Weyl radius and dps-70\\
source-to-$F$ interval weights&NOT PROMOTED&finite census only\\
cofinal quantifier&OPEN&no extrapolation from $k\le16$\\
no RH claim&PASS&external bridges untouched\\
\bottomrule
\end{tabular}
\end{table}

\section{What this does not prove}

It does not prove or disprove a Carleson estimate for infinitely many
selected windows.  It does not prove a source-to-border large sieve.
It does not turn regression into a bound.  It does not infer stability
from a strict SCRAMBLE band that excludes supercritical eigenvalues.
It does not alter any Lean theorem or any claim status.  No RH claim.

\end{document}
